% =============================================================================
%  Sección 4: Conexión con la Tesis — Hacia una Movilidad Sustentable
%             en el Anillo Periférico
% =============================================================================

\section{Conexión con la Tesis: Hacia una Movilidad Sustentable en el Anillo Periférico}

\subsection{El Anillo Periférico como frontera ecológica}
% Análisis de cómo esta vialidad interactúa con los flujos naturales:
%   - El Periférico Sur-Oriente como barrera entre Xochimilco y la ciudad.
%   - Interrupción de corredores bióticos y flujos hídricos superficiales.
%   - Efectos de borde: fragmentación de hábitat, efecto isla, ruido y
%     contaminación lumínica sobre las zonas de anidación.
% Mapeo conceptual de los puntos críticos de intersección
% vialidad--humedal a lo largo del corredor.

\subsection{Replanteando el Corredor de Transporte Público}
% Propuestas de mitigación para integrar la movilidad masiva con la
% conservación de los humedales:
%   - \termino{Infraestructura verde}: camellones arbolados, taludes
%     permeables y jardines de lluvia.
%   - Sistemas de Drenaje Urbano Sostenible (SUDS): pavimento permeable
%     y pozos de infiltración a lo largo del corredor BRT.
%   - Captación y almacenamiento pluvial para recarga de acuíferos.
%   - Pasos de fauna subterráneos cada 500\,m para conectividad
%     entre parches de humedal.
%   - Estaciones del corredor como \termino{nodos verdes}: techos
%     vegetados, paneles solares y compostaje.
% Demostrar que un proyecto de movilidad masiva puede integrarse
% a la conservación si se diseña bajo principios de resiliencia urbana.
